\documentclass[11pt,a4paper]{article}

\usepackage[utf8]{inputenc}
\usepackage[T1]{fontenc}
\usepackage[spanish,es-tabla]{babel}
\usepackage{amsmath,amssymb,amsthm,mathtools}
\usepackage{booktabs}
\usepackage{enumitem}
\usepackage{geometry}
\usepackage{setspace}
\usepackage{hyperref}
\usepackage{microtype}

\geometry{margin=1in}
\onehalfspacing

\hypersetup{
  pdftitle={Nota formal: Ponderacion ontologica alpha(lambda)},
  pdfauthor={Johel Padilla-Villanueva},
  pdfkeywords={RECD, alpha, lambda, ponderacion ontologica, Feigenbaum}
}

\theoremstyle{definition}
\newtheorem{definition}{Definici\'on}[section]
\newtheorem{remark}{Observaci\'on}[section]
\newtheorem{example}{Ejemplo}[section]
\theoremstyle{plain}
\newtheorem{proposition}{Proposici\'on}[section]
\newtheorem{lemma}{Lema}[section]
\newtheorem{theorem}{Teorema}[section]
\newtheorem{corollary}{Corolario}[section]
\newtheorem{conjecture}{Conjetura}[section]
\theoremstyle{remark}
\newtheorem{scholium}{Scholium}[section]
\newtheorem{axiom}{Axioma}[section]

\newcommand{\Res}{\mathrm{Res}}
\newcommand{\RECD}{\mathrm{RECD}}
\newcommand{\TC}{\mathrm{TC}}
\newcommand{\Syn}{\mathrm{Syn}}
\newcommand{\Surp}{\mathrm{Surp}}
\newcommand{\excess}{\mathrm{excess3}}
\newcommand{\Rpair}{\mathcal{R}_{\mathrm{pair}}}
\newcommand{\Jwin}{\mathcal{J}}

\title{\textbf{Nota formal}\\[0.4em]
\large Estatus de la ponderaci\'on $\alpha(\lambda)$:\\
modulación ontológica, principios y l\'imites\\[0.6em]
\normalsize (Punto 3 del roadmap te\'orico; asume Notas 1--2)}

\author{
Johel Padilla-Villanueva\\
\textit{Department of Environmental Health, University of Puerto Rico}\\
\texttt{johelpadilla@gmail.com}
}

\date{2026-08-01 \quad$\cdot$\quad v0.1}

\begin{document}
\maketitle

\begin{abstract}
Esta nota fija el estatus formal de la \emph{ponderaci\'on ontológica}
$\alpha_\ell(\lambda)$ en el RECD anidado. Se separan tres objetos
($\lambda$, $\alpha_\ell$, $f_\ell$), se define con precisi\'on qu\'e
significa ``modulación ontológica'' (y qu\'e \emph{no} significa), se
axiomatiza el motor de pesos, y se clasifica cada pieza como
\textbf{principio de dise\~no}, \textbf{derivable}, o \textbf{par\'ametro
libre preespecificado}. Se exploran cuatro v\'ias de justificación m\'as
profunda (monoton\'ia de profundidad, m\'axima entrop\'ia bajo masa de
claims, estabilidad de conjunciones, consistencia con renormalización
Feigenbaum) y se delimita lo que hoy \emph{no} se puede pretender
demostrar. El resultado: $\alpha(\lambda)$ deja de ser una niebla
retórica y pasa a ser un objeto con axiomas, grados de libertad
explícitos y un programa de derivación parcial.
\end{abstract}

\tableofcontents
\vspace{1em}

% ============================================================
\section{Canon heredado y pregunta del punto 3}
\label{sec:q}

\begin{scholium}
Valen Notas 1--2: tres planos; L3 fuerte $:=$ $\Res_{\mathrm{pair}}>0$;
$\excess$ proxy; anidamiento Tipo~II; M-I-1.
\end{scholium}

El reloj anidado es
\begin{equation}
\label{eq:drecd}
\Delta\RECD(t)
=\sum_{\ell=1}^{3}\alpha_\ell\bigl(\lambda(t)\bigr)\,
\widehat{M}_\ell(t),
\end{equation}
con $\widehat{M}_1=\Phi_1$, $\widehat{M}_2=\Phi_2$,
$\widehat{M}_3\in\{\excess,\Res_{\mathrm{pair}},\Phi_3\}$.

\textbf{Pregunta del punto~3:}
¿qu\'e es exactamente $\alpha(\lambda)$? ¿Par\'ametro fenomenológico
ajustable, o pieza con justificación estructural?

% ============================================================
\section{Tres objetos que no se colapsan}
\label{sec:three}

\begin{definition}[Cue de r\'egimen $\lambda$]
\label{def:lambda}
Una funci\'on $\lambda(t)\ge 0$ (o $\lambda\in[0,1]$ tras normalizar)
construida a partir de observables de r\'egimen---p.ej.\ $|\tau_s|$,
espectro, RQA, par\'ametro de control $r$---con el
\textbf{design goal}
\begin{equation}
\label{eq:goal}
\mathbb{E}\bigl[\lambda\mid\mathcal{R}_{\mathrm{chaos}}\bigr]
>
\mathbb{E}\bigl[\lambda\mid\mathcal{R}_{\mathrm{pre}}\bigr].
\end{equation}
$\lambda$ \emph{no} es un claim de conjunción; es una pista de en qu\'e
r\'egimen opera el sistema.
\end{definition}

\begin{definition}[Agenda de pesos $\alpha_\ell(\lambda)$]
\label{def:alpha}
Una terna de funciones
$\alpha_\ell:[0,\infty)\to(0,\infty)$, $\ell=1,2,3$, que convierten el
cue en coeficientes del incremento \eqref{eq:drecd}.
\end{definition}

\begin{definition}[Masa ontológica de nivel / share $f_\ell$]
\label{def:f}
Siempre que el denominador sea positivo,
\begin{equation}
\label{eq:f}
f_\ell(t)
:=\frac{\alpha_\ell\bigl(\lambda(t)\bigr)\,\widehat{M}_\ell(t)}
{\sum_{k=1}^{3}\alpha_k\bigl(\lambda(t)\bigr)\,\widehat{M}_k(t)}
\in[0,1].
\end{equation}
La terna $(f_1,f_2,f_3)$ es la \textbf{composici\'on} del tick RECD en $t$.
\end{definition}

\begin{proposition}[Independencia de planos]
\label{prop:indep}
\begin{enumerate}[label=(\roman*)]
  \item Los claims $C_\ell$ y medidas $M_\ell$ (Notas 1--2) no dependen de
  $\alpha$.
  \item $\lambda$ puede calcularse sin RECD (p.ej.\ solo $\tau_s$).
  \item $f_\ell$ mezcla \emph{evidencia de claims} ($\widehat{M}_\ell$)
  con \emph{pol\'itica de pesos} ($\alpha_\ell$): no es un observable
  puro del joint.
\end{enumerate}
\end{proposition}

\begin{scholium}[Consecuencia editorial]
La abundancia cruda de L3 ($\mathbb{E}[\excess]$, highL3,
$\Res_{\mathrm{pair}}$) es un resultado de Nivel~3 \emph{sin} $\lambda$.
La share $f_3$ es un resultado \emph{con} motor de pesos. Confundirlos
es el error m\'as frecuente al leer validaciones.
\end{scholium}

% ============================================================
\section{Qu\'e significa ``modulación ontológica''}
\label{sec:meaning}

\begin{definition}[Modulación ontológica --- sentido operativo]
\label{def:mod}
Llamamos \textbf{modulación ontológica} al mapa
\begin{equation}
\bigl(\widehat{M}_1,\widehat{M}_2,\widehat{M}_3;\lambda\bigr)
\;\longmapsto\;
\bigl(f_1,f_2,f_3;\,\Delta\RECD\bigr)
\end{equation}
inducido por $\alpha_\ell(\lambda)$. El adjetivo ``ontológica'' denota
que el reloj no trata todos los eventos ordinales como equivalentes en
\emph{masa de tick}: la profundidad del claim modula cu\'anto avanza el
tiempo extramental discreto.
\end{definition}

\begin{scholium}[Lo que \emph{no} significa]
\label{sch:not}
\begin{enumerate}[label=(\roman*)]
  \item No significa que el caos ``cree el ser'' ni que $\alpha$ mida el
  \emph{actus essendi} (M-I-1).
  \item No significa que $\lambda$ sea una variable ontológica del
  mundo; es un cue epist\'emico/operativo de r\'egimen.
  \item No significa que $f_3$ alto demuestre sinergia irreducible: $f_3$
  puede subir por $\alpha_3(\lambda)$ aunque $\widehat{M}_3$ sea un proxy
  d\'ebil (Nota~2).
  \item No introduce un segundo reloj: es el \emph{mismo} RECD con
  coeficientes r\'egimen-dependientes (absorción del gate legacy
  $g(\tau_s)$, $\delta^{-k}$).
\end{enumerate}
\end{scholium}

\begin{remark}[Intuición controlada]
Nivel~1: eventos frecuentes, delgados (coincidencia).
Nivel~2: memoria relacional.
Nivel~3: residuo fuera de $\Rpair$.
La hip\'otesis de ponderación: cerca de regímenes de alta volatilidad /
caos desarrollado, la \emph{composici\'on} del tick se desplaza hacia
capas m\'as profundas ($f_3\uparrow$ relativo a $f_1$).
\end{remark}

% ============================================================
\section{Axiomas del motor de pesos}
\label{sec:axioms}

Un motor $(\lambda,\alpha)$ es \textbf{admisible} si cumple:

\begin{axiom}[Positividad]
\label{ax:pos}
$\alpha_\ell(\lambda)>0$ para todo $\ell$ y todo $\lambda\ge 0$ en el
dominio. Ning\'un nivel se apaga por completo (evita singularidades en
$f_\ell$ y p\'erdida de identificabilidad).
\end{axiom}

\begin{axiom}[Monoton\'ia de profundidad]
\label{ax:mono}
Existen funciones mon\'otonas (no necesariamente estrictas) tales que,
al crecer $\lambda$ en el r\'egimen de inter\'es,
\begin{equation}
\alpha_1(\lambda)\downarrow,\qquad
\alpha_2(\lambda)\uparrow\text{ o se estabiliza},\qquad
\alpha_3(\lambda)\uparrow.
\end{equation}
Es decir: el cue de caos/alta volatilidad \emph{no reduce} la masa
asignable a claims m\'as profundos.
\end{axiom}

\begin{axiom}[Separación cue / claim]
\label{ax:sep}
$\lambda$ no se define como funci\'on exclusiva de $\widehat{M}_3$ con
peso dominante (circularidad: reponderar L3 con L3). Si se usa un cue
lento de L3, su peso en $\lambda$ est\'a acotado
($w_{\mathrm{L3}}\le 0.15$ en el contrato multi-cue del paper de
fundamentos).
\end{axiom}

\begin{axiom}[Preespecificación]
\label{ax:pre}
La forma de $\alpha_\ell$ y la construcci\'on de $\lambda$ se fijan
\emph{a priori} (o en segmento de calibración etiquetado separado del
test). No se optimizan sobre el dataset de descubrimiento de L3.
\end{axiom}

\begin{axiom}[Design goal del cue]
\label{ax:goal}
$\lambda$ satisface \eqref{eq:goal} en el sistema de inter\'es (tras
polaridad $s$ o multi-cue). Si no puede verificarse, se reporta RECD
\emph{sin} reponderar ($\alpha_\ell\equiv 1$) o se declara el fallo del
cue.
\end{axiom}

\begin{proposition}[Template de referencia es admisible]
\label{prop:template}
El template del paper de fundamentos
\begin{equation}
\label{eq:template}
\begin{aligned}
\alpha_1(\lambda)&=\alpha_{10}\,e^{-\beta_1\lambda},\\
\alpha_2(\lambda)&=\alpha_{20}\,(1+\gamma_2\lambda),\\
\alpha_3(\lambda)&=\alpha_{30}\,(1+\gamma_3\lambda+\delta_3\lambda^2),
\end{aligned}
\end{equation}
con $\alpha_{\ell 0}>0$ y $\beta_1,\gamma_2,\gamma_3,\delta_3\ge 0$,
cumple Axiomas~\ref{ax:pos}--\ref{ax:mono}. Los Axiomas~\ref{ax:sep}--\ref{ax:goal}
dependen de c\'omo se construya $\lambda$, no de la forma \eqref{eq:template}.
\end{proposition}

\begin{proof}
Derivadas: $\alpha_1'=-\beta_1\alpha_1\le 0$;
$\alpha_2'=\alpha_{20}\gamma_2\ge 0$;
$\alpha_3'=\alpha_{30}(\gamma_3+2\delta_3\lambda)\ge 0$; positividad
obvia.
\end{proof}

% ============================================================
\section{Clasificaci\'on de estatus: qu\'e es cada pieza}
\label{sec:status}

\begin{definition}[Tres estatus]
\label{def:status}
\begin{enumerate}[label=(\roman*)]
  \item \textbf{Derivable:} se obtiene de un principio enunciado +
  hip\'otesis, sin knobs libres esenciales.
  \item \textbf{Principio de dise\~no:} se postula por coherencia con el
  paradigma (axiomas), no se ``mide'' en la naturaleza.
  \item \textbf{Par\'ametro libre preespecificado:} forma o n\'umero fijado
  a priori por contrato metodológico; puede calibrarse en holdout, no en
  el test de L3.
\end{enumerate}
\end{definition}

\begin{center}
\begin{tabular}{@{}p{0.34\textwidth}p{0.22\textwidth}p{0.36\textwidth}@{}}
\toprule
Pieza & Estatus actual & Comentario \\
\midrule
Claims $C_\ell$, $M_\ell$, $\Res_{\mathrm{pair}}$ &
Derivable / objeto & Notas 1--2 \\
Design goal \eqref{eq:goal} &
Principio de dise\~no & Motor debe elevar L3 donde se espera \\
Axiomas \ref{ax:pos}--\ref{ax:pre} &
Principio de dise\~no & Admisibilidad del motor \\
Forma cualitativa $\alpha_1\!\downarrow,\alpha_3\!\uparrow$ &
Principio de dise\~no & Monoton\'ia de profundidad \\
Template exponencial/polinomial \eqref{eq:template} &
Par\'ametro libre preesp. & Una familia admisible, no \'unica \\
$(\beta_1,\gamma_2,\gamma_3,\delta_3)$ &
Par\'ametro libre preesp. & Defaults de referencia; no teorema \\
Polaridad $s\in\{+1,-1\}$ &
Calibración de sistema & Diagn\'ostico en etiquetas pre/chaos \\
Factor $(1+\gamma\log\delta)$ &
Ancla de renormalización & Continuidad con gate Feigenbaum legacy \\
$\lambda_\tau$ solo con $|\tau_s|$ &
Proxy de cue (fr\'agil) & Puede invertir; exige $s$ o multi-cue \\
$\lambda(r)$ con $r$ conocido &
Cue casi ideal & Separa r\'egimen de ruido de proxy \\
$f_\ell$ &
Constructo r\'egimen-aware & Hereda validez de $\lambda$ \\
\bottomrule
\end{tabular}
\end{center}

\begin{theorem}[Estatus global de $\alpha(\lambda)$]
\label{thm:status}
En el estado actual del paradigma:
\begin{enumerate}[label=(\roman*)]
  \item $\alpha(\lambda)$ \textbf{no} es un observable del joint
  simb\'olico ni un corolario de $\Res_{\mathrm{pair}}$.
  \item $\alpha(\lambda)$ \textbf{es} un \emph{principio de dise\~no
  axiomatizado} (Axiomas~\ref{ax:pos}--\ref{ax:goal}) que implementa la
  tesis: \emph{el tiempo discreto del RECD pondera m\'as los claims m\'as
  profundos cuando el cue de r\'egimen se\~nala caos/alta volatilidad}.
  \item La \emph{familia concreta} \eqref{eq:template} es un
  \textbf{representante preespecificado} del cono de motores admisibles,
  no la \'unica funci\'on compatible con los axiomas.
  \item Por tanto, llamar a $\alpha$ ``modulación ontológica'' es
  leg\'itimo en el sentido operativo de Def.~\ref{def:mod}, e ileg\'itimo
  si se lee como derivación metaf\'isica o como ley natural \'unica.
\end{enumerate}
\end{theorem}

% ============================================================
\section{Cuatro v\'ias hacia justificación m\'as profunda}
\label{sec:vias}

Ninguna est\'a cerrada como teorema \'unico; cada una acota grados de
libertad.

\subsection{V\'ia A --- Monoton\'ia de profundidad (ya axioma)}

Es la m\'inima: Axioma~\ref{ax:mono}. No fija la forma funcional, solo el
cono
\[
\mathcal{A}
=\bigl\{\alpha:\,
\alpha_1'\le 0,\;
\alpha_2'\ge 0,\;
\alpha_3'\ge 0,\;
\alpha_\ell>0\bigr\}.
\]

\begin{proposition}
\label{prop:cone}
Si $\alpha,\tilde\alpha\in\mathcal{A}$ y $\lambda$ cumple
\eqref{eq:goal}, entonces ambas elevan (en media) la share esperada de
L3 al pasar pre$\to$chaos \emph{condicional} a $\widehat{M}_\ell$ fijos en
ley. La magnitud del elevamiento es familia-dependiente.
\end{proposition}

\subsection{V\'ia B --- M\'axima entrop\'ia / Gibbs sobre profundidad}

\begin{definition}[Pesos de Gibbs de profundidad]
\label{def:gibbs}
Sea $u(\ell)$ una utilidad de profundidad ($u(1)<u(2)<u(3)$) y
$\beta(\lambda)\ge 0$ una ``inversa de temperatura'' creciente en
$\lambda$. Def\'inase
\begin{equation}
\label{eq:gibbs}
\alpha_\ell^{\mathrm{G}}(\lambda)
\propto
\exp\bigl(\beta(\lambda)\,u(\ell)\bigr).
\end{equation}
\end{definition}

\begin{proposition}
\label{prop:gibbs_adm}
Si $\beta'(\lambda)\ge 0$ y $u(1)<u(2)<u(3)$, entonces
$\alpha^{\mathrm{G}}$ satisface la monoton\'ia relativa
$\alpha_3/\alpha_1=\exp\bigl(\beta(u(3)-u(1))\bigr)$ creciente en
$\lambda$, y puede reescalarse para cumplir Axioma~\ref{ax:pos}.
Es un motor admisible \emph{derivado} de un principio variacional
est\'andar (m\'axima entrop\'ia bajo media de utilidad de profundidad).
\end{proposition}

\begin{remark}
La V\'ia B \textbf{deriva la forma cualitativa} (más peso a $\ell$ alto
cuando $\beta$ sube), pero \textbf{no deriva} $u(\ell)$ ni
$\beta(\lambda)$ desde los datos de $S_t$ sin hip\'otesis extra. Reduce
knobs: de cuatro exponentes a una funci\'on $\beta$ y una utilidad
discreta de tres valores.
\end{remark}

\begin{conjecture}[Puente template $\leftrightarrow$ Gibbs]
\label{conj:gibbs}
Existen $u$, $\beta$ y un reescalado local tales que
\eqref{eq:template} aproxima $\alpha^{\mathrm{G}}$ en un intervalo
compacto de $\lambda$ con error uniformemente acotado (p.ej.\ en
$\lambda\in[0,2]$). \'Util para reinterpretar defaults, no para
``demostrar'' el template.
\end{conjecture}

\subsection{V\'ia C --- Estabilidad de conjunciones}

Idea: $\alpha_\ell$ deber\'ia crecer con la \emph{estabilidad esperada}
del claim $\ell$ bajo perturbaciones que preservan el r\'egimen.

\begin{definition}[Estabilidad emp\'irica de nivel]
\label{def:stab}
Sea $\mathcal{N}$ una clase de perturbaciones (ruido de medida mon\'otono,
jitter de ventana, surrogados que preservan m\'argenes). Def\'inase
\begin{equation}
\mathrm{Stab}_\ell
:=\mathbb{E}\bigl[
\mathrm{corr}\bigl(\widehat{M}_\ell,\,\widehat{M}_\ell\circ\mathcal{N}\bigr)
\bigr]
\end{equation}
o una tasa de signo/orden preservado.
\end{definition}

\begin{conjecture}[Pesos por estabilidad relativa]
\label{conj:stab}
En regímenes caóticos, $\mathrm{Stab}_3/\mathrm{Stab}_1$ no colapsa a 0
(el residuo fuerte es m\'as fr\'agil en estimación pero m\'as
``informativo'' por evento). Una regla
$\alpha_\ell\propto g(\mathrm{Stab}_\ell,\lambda)$ con $g$ creciente en
ambos puede \emph{calibrar} $\alpha$ sin optimizar sobre el label de
transici\'on. \textbf{Estado:} programa abierto; riesgo de
circularidad si $\mathcal{N}$ se elige ad hoc.
\end{conjecture}

\subsection{V\'ia D --- Consistencia ordinal con Feigenbaum}

\begin{definition}[Factor de escala Feigenbaum en el cue]
\label{def:feig}
Con $\delta\approx 4.6692016091$,
\begin{equation}
\kappa_\delta:=1+\gamma\log\delta,\qquad\gamma\ge 0,
\end{equation}
y $\lambda$ multiplica o incorpora $\kappa_\delta$ (template
$\lambda_\tau$ del paper). En la rama legacy, el paso de tiempo en la
cascada usaba $\delta^{-k}$; aqu\'i esa contracción se reinterpreta como
\emph{aumento de masa por profundidad} cuando el cue entra en r\'egimen
ca\'otico, no como un segundo reloj.
\end{definition}

\begin{axiom}[No contradicción con renormalización]
\label{ax:feig}
Si el sistema atraviesa doblamientos de periodo con acumulaci\'on en
$r_\infty$, el motor $(\lambda,\alpha)$ no debe \emph{invertir} la
predicci\'on de masa L3 al cruzar $r_\infty$ cuando $\lambda$ est\'a
bien calibrado (design goal). El factor $\kappa_\delta$ es una ancla de
continuidad con la teor\'ia de renormalización, no una demostración de
que $\delta$ sea la causa eficiente del tick.
\end{axiom}

\begin{conjecture}[$\alpha$ a lo largo de la cascada; puente a pto.~4]
\label{conj:cascade}
En mapas unimodales acoplados, con $\lambda=\lambda(r)$ mon\'otono en el
par\'ametro de control:
\begin{enumerate}[label=(\roman*)]
  \item $f_3(r)$ es no decreciente al pasar por la secuencia de
  doblamientos hacia $r_\infty$ y el caos desarrollado, en media;
  \item el salto cualitativo de composici\'on $(f_1,f_2,f_3)$ cerca de
  $r_\infty$ no se reduce a un rescaling de $|\tau_s|$;
  \item $\alpha_3(\lambda(r))/\alpha_1(\lambda(r))$ rastrea una
  funci\'on mon\'otona de la ``profundidad de cascada'' (p.ej.\ $k$ en
  periodo $2^k$ antes de $r_\infty$).
\end{enumerate}
\end{conjecture}

% ============================================================
\section{Construcción de $\lambda$: estatus y jerarqu\'ia}
\label{sec:lambda_build}

\begin{definition}[Jerarqu\'ia de fallback (canon operativo)]
\label{def:fallback}
\begin{enumerate}[label=(\arabic*)]
  \item $\lambda=\lambda(r)$ mon\'otono si hay control ground-truth.
  \item $\lambda_{\mathrm{multi}}$ (consenso $\tau$ + espectro + RQA;
  $w_{\mathrm{L3}}\le 0.15$ opcional).
  \item $\lambda_\tau$ con polaridad $s$ diagnosticada en holdout
  etiquetado, solo si se verifica \eqref{eq:goal}.
  \item $\alpha_\ell\equiv 1$ (RECD sin reponderar) si no hay cue
  fiable.
\end{enumerate}
\end{definition}

\begin{theorem}[Fragilidad del cue $|\tau_s|$ puro]
\label{thm:fragile}
No es teorema que $|\tau_s|$ sea mon\'otono con el caos en toda clase de
sistemas. Existen acoplamientos (p.ej.\ log\'isticos del paper de
fundamentos) donde $\overline{|\tau_s|}_{\mathrm{pre}}>\overline{|\tau_s|}_{\mathrm{chaos}}$.
Por tanto $s$ o multi-cue son \textbf{parte del estatus formal} de
$\lambda$, no detalles de implementaci\'on.
\end{theorem}

\begin{proposition}[Polaridad como isometr\'ia de r\'egimen]
\label{prop:polarity}
El mapa $s\cdot g(|\tau_s|)$ con $s\in\{+1,-1\}$ es la corrección
m\'inima que restaura \eqref{eq:goal} cuando el orden pre/chaos de
$|\tau_s|$ se invierte. No altera los claims $M_\ell$.
\end{proposition}

% ============================================================
\section{Qu\'e se conserva y qu\'e se pierde al ponderar}
\label{sec:conserve}

\begin{proposition}[Invariantes bajo $\alpha$]
\label{prop:inv_alpha}
\begin{enumerate}[label=(\roman*)]
  \item $\widehat{M}_\ell(t)$ y $\Res_{\mathrm{pair}}(\Jwin_t)$ son
  invariantes bajo cambio de $\alpha$ (Plano claim/estimador).
  \item El signo de $\Delta\tau_s$ y las bandas de $\tau_s$ no dependen
  de $\alpha$.
  \item La ordenaci\'on de dos trayectorias por $T_{\RECD}$ \emph{puede}
  cambiar al cambiar $\alpha$ (la ponderación altera el reloj
  acumulado).
\end{enumerate}
\end{proposition}

\begin{proposition}[Identificabilidad]
\label{prop:id}
Si $\widehat{M}_\ell$ son linealmente dependientes en una ventana, $f_\ell$
no identifica $\alpha$ de forma \'unica. En particular, reportar solo
$f_3$ sin $(\widehat{M}_\ell,\lambda,\alpha)$ es subidentificado.
\end{proposition}

\begin{scholium}[Norma de reporte]
Todo resultado que invoque ``masa ontológica'' debe publicar:
(1)~$\widehat{M}_\ell$ crudos; (2)~construcci\'on de $\lambda$ y
verificación de \eqref{eq:goal}; (3)~$\alpha_\ell$ usados;
(4)~$f_\ell$ como derivado. Sin (1)--(2), $f_3$ no es interpretable.
\end{scholium}

% ============================================================
\section{Respuesta n\'itida a la pregunta del roadmap}
\label{sec:answer}

\begin{enumerate}[label=\textbf{Q\arabic*.}, leftmargin=*]
  \item \textbf{¿Qu\'e significa ``modulación ontológica''?}\\
  El mapa operativo
  $(\widehat{M},\lambda)\mapsto(f,\Delta\RECD)$: la profundidad del claim
  modula la masa del tick del reloj extramental discreto
  (Def.~\ref{def:mod}). No es metaf\'isica del ser
  (Scholium~\ref{sch:not}).

  \item \textbf{¿Es $\alpha(\lambda)$ fenomenológico o estructural?}\\
  \textbf{Ambos, en capas:}
  \begin{itemize}
    \item \emph{Estructural / de principio:} axiomas de admisibilidad y
    monoton\'ia de profundidad (el reloj \emph{debe} poder elevar L3
    cuando el r\'egimen lo exige).
    \item \emph{Fenomenológico / preespecificado:} la forma concreta
    \eqref{eq:template} y sus coeficientes num\'ericos.
    \item \emph{Parcialmente derivable:} V\'ia B (Gibbs/m\'axent sobre
    profundidad) y ancla Feigenbaum $\kappa_\delta$ reducen libertad
    sin cerrar un teorema \'unico.
  \end{itemize}

  \item \textbf{¿Puede derivarse de un principio m\'as profundo?}\\
  Hoy: \textbf{derivación parcial} (cono $\mathcal{A}$, Gibbs,
  design goal, no-contradicción Feigenbaum). No hay a\'un un principio
  \'unico que fije $\alpha_1,\alpha_2,\alpha_3$ sin knobs. Las Conjs.\
  \ref{conj:gibbs}--\ref{conj:cascade} marcan el programa.

  \item \textbf{¿C\'omo cambia $\alpha(\lambda)$ en la ruta de
  Feigenbaum?}\\
  Si $\lambda=\lambda(r)$ mon\'otono, $\alpha_3/\alpha_1$ crece hacia y
  m\'as all\'a de $r_\infty$ por construcci\'on del template admisible.
  Que $f_3$ y $\Res_{\mathrm{pair}}$ hagan lo mismo de forma no trivial
  es la Conj.~\ref{conj:cascade} (punto~4).
\end{enumerate}

% ============================================================
\section{Canon a\~nadido (C15--C22)}
\label{sec:canon}

\begin{enumerate}[label=\textbf{C\arabic*.}]
  \setcounter{enumi}{14}
  \item Separar siempre $\lambda$, $\alpha_\ell$, $f_\ell$.
  \item Modulación ontológica $:=$ Def.~\ref{def:mod} (operativa).
  \item Motor admisible $:=$ Axiomas~\ref{ax:pos}--\ref{ax:goal}.
  \item $\alpha(\lambda)$ es principio de dise\~no axiomatizado + familia
  preespecificada; no es observable del joint.
  \item Abundancia L3 (cruda) $\ne$ share $f_3$ (reponderada).
  \item Fallback de $\lambda$: $r$ $\to$ multi-cue $\to$ $\lambda_\tau$
  calibrado $\to$ $\alpha\equiv 1$.
  \item Polaridad $s$ es parte del estatus formal del cue $|\tau_s|$.
  \item Norma de reporte: $(\widehat{M}_\ell,\lambda,\alpha,f_\ell)$.
\end{enumerate}

\subsection{Implementaci\'on emp\'irica V\'ia B (v0.2)}
\label{sec:gibbs_emp}

Software: \texttt{nested-recd} $\ge 0.2.2$ expone
\texttt{alpha\_weights\_gibbs} y
\texttt{alpha\_compare\_template\_gibbs}.
Script: \texttt{scripts/gibbs\_alpha\_compare.py}.
Defaults: $u=(0,1,2)$, $\beta(\lambda)=\kappa\lambda$ con $\kappa=1.5$.

\paragraph{Puente de forma (Conj.~\ref{conj:gibbs}).}
Sobre $\lambda\in[0,2]$, familias L1-normalizadas:
$\max L_\infty\approx 0.163$, $\max L_1\approx 0.326$,
$\mathrm{mean}\,L_1\approx 0.215$. Aproximaci\'on de forma
\emph{acotada}, no identidad.

\paragraph{Share $f_3$ en la cascada (misma masa $\Phi$).}
Estaciones $r\in\{3.20,\ldots,3.85\}$, $N_{\mathrm{real}}=4$:
$\mathrm{corr}_r(f_3^{\mathrm{T}},f_3^{\mathrm{G}})\approx 0.997$
(ambos motores empujan el share de L3 en el mismo sentido cualitativo);
$\mathrm{corr}_r(f_3^{\mathrm{T}},\mathrm{excess3})\approx 0.08$
--- refuerzo emp\'irico de \textbf{abundancia $\ne$ share}.
En $r\le r_\infty$ el proxy $\lambda(r)=0$ deja ambos motores en el
baseline; en caos desarrollado ($r=3.85$) Gibbs eleva levemente
$f_3$ ($0.926$ vs $0.884$ template).

\begin{scholium}
La V\'ia B est\'a implementada y es intercambiable \emph{cualitativamente}
con el template en esta rejilla; no sustituye la preespecificaci\'on
ni convierte $\alpha$ en observable del joint.
\end{scholium}

\vspace{1.5em}
\noindent\rule{\textwidth}{0.4pt}\\[0.5em]
\begin{center}
\em Fin de la Nota formal punto 3 v0.2
(V\'ia B implementada + puente emp\'irico).\\
Artefactos: \texttt{results/gibbs\_alpha/},
\texttt{figures/gibbs\_*.png}; software \texttt{nested-recd} 0.2.2.
\end{center}

\begin{thebibliography}{99}
\bibitem{Padilla2026note1}
J.~Padilla-Villanueva.
\newblock Nota formal 1: Conjunciones ordinales\ldots
\newblock Fundaci\'on RECD, 2026.
\bibitem{Padilla2026note2}
J.~Padilla-Villanueva.
\newblock Nota formal 2: Sinergia irreducible\ldots
\newblock Fundaci\'on RECD, 2026.
\bibitem{Padilla2026framework}
J.~Padilla-Villanueva.
\newblock Systemic Tau and the RECD Framework.
\newblock Manuscrito de fundamentos, 2026.
\end{thebibliography}

\end{document}
